% ============================================================
% Capítulo 7 — Language Network: conceptos + código
% ============================================================
\chapter{Red de lenguaje: BERT, embeddings, BiGRU y pack-padded}
\label{cap:language-network}

Este capítulo explica los conceptos de NLP detrás de la rama
textual del modelo y luego muestra el código línea por línea.

\section{¿Qué es BERT y por qué sólo su tokenizer?}

\subsection{BERT}

BERT (\emph{Bidirectional Encoder Representations from
Transformers})~\cite{devlin2018bert} es un modelo de lenguaje
preentrenado en un corpus enorme de texto. Tiene dos partes:

\begin{enumerate}
  \item \textbf{Tokenizer}: convierte texto a IDs numéricos usando
        un vocabulario fijo. Es lo único que usamos aquí.
  \item \textbf{Transformer encoder}: 12 capas de atención que
        producen embeddings contextuales. \textbf{No lo usamos}.
\end{enumerate}

\subsection{¿Por qué no usar BERT entero?}

BERT completo tiene \textasciitilde 110M parámetros. BiGRU sólo
\textasciitilde 4.7M. Para textos cortos (descripciones de 30
palabras) BiGRU captura suficiente contexto. Esto reduce:

\begin{itemize}
  \item Coste computacional (\textasciitilde 20$\times$ menos).
  \item VRAM (cabe en GPU modestas).
  \item Tiempo de entrenamiento.
\end{itemize}

\subsection{Cómo funciona el tokenizer de BERT}

BERT usa \emph{WordPiece tokenization}: el vocabulario está
formado por sub-palabras (\emph{tokens}) frecuentes más letras
individuales. Ejemplo:

\begin{verbatim}
Entrada:  "the man is wearing a black jacket"
Tokens:   ['the', 'man', 'is', 'wearing', 'a', 'black', 'jacket']
IDs:      [1996, 2158, 2003, 7078, 1037, 2304, 7505]
Especial: [CLS] = 101, [SEP] = 102
\end{verbatim}

\noindent\textbf{Vocabulario}: 30\,522 tokens base. Más 1 (índice
0 reservado para padding) $\approx$ 30\,609 únicos. Más 1 para
\texttt{[PAD]} = 29\,610 (lo que usa este proyecto).

\section{¿Qué es un embedding?}

Un \emph{embedding} es una representación densa (vector de
números reales) de un token. En este proyecto:

\begin{itemize}
  \item Cada token del vocabulario $\to$ vector de 512 números
        reales.
  \item Total de parámetros: $29\,610 \times 512 = 15,16$M.
\end{itemize}

\noindent\textbf{Inicialización}: $\mathcal{N}(0, 1)$ por defecto.
Durante el entrenamiento, estos vectores se ajustan para que
tokens con significado similar tengan vectores cercanos.

\section{¿Qué es un padding?}

Los textos del batch tienen longitudes distintas. Para meterlos
en un tensor uniforme se \emph{padean} a una longitud fija (100
en este proyecto):

\begin{verbatim}
"a man in red"     -> [101, 1037, 2158, 1999, 2417, 102] (6 tokens)
"a tall person"    -> [101, 1037, 4206, 2718, 102] (5 tokens)
\end{verbatim}

\noindent Ambos se convierten a longitud 100 añadiendo ceros
(\texttt{[PAD]} = 0) al final. El modelo \textbf{ignora} el padding
porque:

\begin{enumerate}
  \item El índice 0 nunca se usa como token real (es
        \texttt{[PAD]}).
  \item El embedding de índice 0 está fijo en $(0, 0, \ldots, 0)$.
  \item La GRU recibe \texttt{orig\_token\_length} para saber
        dónde termina cada secuencia.
\end{enumerate}

\section{¿Qué es una GRU?}

Una \emph{Gated Recurrent Unit}~\cite{cho2014gru} es una versión
simplificada de LSTM. Procesa una secuencia token por token y
mantiene un \emph{estado oculto} $h_t$:

\begin{align}
  z_t &= \sigma(W_z \cdot [h_{t-1}, x_t])  &\text{(update gate)} \\
  r_t &= \sigma(W_r \cdot [h_{t-1}, x_t])  &\text{(reset gate)} \\
  \tilde{h}_t &= \tanh(W \cdot [r_t \odot h_{t-1}, x_t]) \\
  h_t &= (1 - z_t) \odot h_{t-1} + z_t \odot \tilde{h}_t
\end{align}

\noindent\textbf{Bidireccional} significa que se corren dos GRU en
paralelo: una hacia adelante y otra hacia atrás. Las salidas se
\emph{promedian} para tener en cuenta contexto en ambas
direcciones.

\noindent\textbf{Ventajas} sobre LSTM:
\begin{itemize}
  \item Menos parámetros (2 compuertas vs 3).
  \item Más rápido en secuencias cortas.
\end{itemize}

\section{¿Qué es pack-padded-sequence?}

\subsection{El problema}

Si procesamos un batch con secuencias de longitudes [10, 30, 5,
20] todas padeadas a 100, la GRU hace trabajo inútil en los
\texttt{[PAD]}. Además los \texttt{[PAD]} afectan el estado
oculto: como $h_{100} = f(h_{99}, 0)$, el estado final está
"contaminado" por 70 pasos de padding.

\subsection{La solución}

\texttt{pack\_padded\_sequence} de PyTorch:

\begin{enumerate}
  \item \textbf{Ordena} las secuencias por longitud descendente.
  \item \textbf{Elimina} los pasos de padding antes de pasar a
        la RNN.
  \item La RNN procesa sólo los tokens reales.
  \item \texttt{pad\_packed\_sequence} restaura el orden original
        con padding.
\end{enumerate}

\noindent\textbf{Requisito}: pasar las longitudes reales en CPU
(no GPU) porque el backend de PyTorch así lo requiere.

\section{Código: \texttt{GRULanguageNetwork}}

\begin{minted}[language=Python, numbers, firstnumber=11, linerange=11-41]{python}
class GRULanguageNetwork(nn.Module):
    """Embedding + Dropout + GRU bidireccional."""

    def __init__(self, config=None):
        super(GRULanguageNetwork, self).__init__()
        if config is None:
            raise ValueError("`config` can not be none")

        # Capa de embedding
        # vocab_size = 29609 (BERT) + 1 (padding) = 29610
        # embedding_dim = 512
        # padding_idx = 0: el vector para índice 0 se mantiene en 0
        self.embedding = nn.Embedding(
            config.vocab_size,        # 29610
            config.embedding_dim,     # 512
            padding_idx=0
        )

        # Dropout tras el embedding (30%)
        self.dropout = nn.Dropout(config.dropout_rate)  # 0.30

        # GRU bidireccional
        # input_size: 512 (dim del embedding)
        # hidden_size: 1024 (feature_length)
        # num_layers: 1
        # bidirectional=True: procesa en ambas direcciones
        # bias=False: ahorra parámetros
        self.gru = nn.GRU(
            config.embedding_dim,     # 512
            config.feature_length,    # 1024
            num_layers=config.num_layers,    # 1
            bidirectional=True,
            bias=False
        )

    def forward(self, text_ids=None, text_length=None):
        """text_ids: (B, 100), text_length: (B,).
           returns: (B, 1024, 100, 1)"""
        # 1. Embedding: (B, 100) -> (B, 100, 512)
        text_embedding = self.embedding(text_ids)
        # 2. Dropout (solo en training mode)
        text_embedding = self.dropout(text_embedding)
        # 3. Procesar con BiGRU (usando pack_padded_sequence)
        feature = self.do_RNN_computation(text_embedding, text_length, self.gru)
        return feature
\end{minted}

\section{Código: \texttt{do\_RNN\_computation}}

\begin{minted}[language=Python, numbers, firstnumber=44, linerange=44-69]{python}
    def do_RNN_computation(self, text_embedding, text_length, gru):
        """Procesa la secuencia con GRU usando pack_padded_sequence."""

        # Asegurar que text_length es 1-D
        text_length = text_length.view(-1)

        # Ordenar índices por longitud DESCENDENTE
        _, sort_index = torch.sort(text_length, dim=0, descending=True)

        # Invertir el orden: nos servirá para "des-ordenar" después
        _, unsort_index = sort_index.sort()

        # Reordenar el batch por longitud descendente
        sortlength_text_embedding = text_embedding[sort_index, :]
        sort_text_length = text_length[sort_index]

        # Empaquetar: ignora el padding en el cómputo de la GRU
        packed_text_embedding = pack_padded_sequence(
            sortlength_text_embedding,
            sort_text_length.cpu(),
            batch_first=True
        )

        # Pasar por la GRU
        packed_feature, hn = gru(packed_text_embedding)

        # Desempaquetar: vuelve a la forma (B, T, *)
        total_length = text_embedding.size(1)
        sort_feature = pad_packed_sequence(
            packed_feature, batch_first=True, total_length=total_length
        )

        # Des-ordenar: volver al orden original del batch
        unsort_feature = sort_feature[0][unsort_index, :]

        # Promediar fwd y bwd:
        # BiGRU output: (B, T, 2*hidden) = (B, 100, 2048)
        # [:, :, :1024] = forward, [:, :, 1024:] = backward
        unsort_feature = (
            unsort_feature[:, :, :int(unsort_feature.size(2) / 2)] +
            unsort_feature[:, :, int(unsort_feature.size(2) / 2):]
        ) / 2

        # Preparar para alimentar el DSC final:
        # permute(0,2,1): (B, 100, 1024) -> (B, 1024, 100)
        # unsqueeze(3): añadir dim al final -> (B, 1024, 100, 1)
        return unsort_feature.permute(0, 2, 1).contiguous().unsqueeze(3)
\end{minted}

\section{Código: \texttt{BERTTokenizer}}

\begin{minted}[language=Python, numbers, firstnumber=10, linerange=10-19]{python}
class BERTTokenizer(object):
    """Adapta el tokenizer de HuggingFace a la interfaz del proyecto."""

    def __init__(self):
        super(BERTTokenizer, self).__init__()

        # Importar transformers
        _, tokenizer_class, pretrained_weights = (
            ppb.BertModel, ppb.BertTokenizer, 'bert-base-uncased'
        )

        # Cargar el tokenizer preentrenado (BERT-base, vocab inglés)
        self.tokenizer = tokenizer_class.from_pretrained(pretrained_weights)

    def __call__(self, text: str = None) -> torch.LongTensor:
        """Codifica texto a IDs de tokens."""
        # .encode() añade automáticamente [CLS] y [SEP]
        tokens = self.tokenizer.encode(text)
        return tokens
\end{minted}

\section{Código: \texttt{pad\_tokens}}

\begin{minted}[language=Python, numbers, firstnumber=86, linerange=86-107]{python}
def pad_tokens(tokens, tokens_length_max):
    """Padea o trunca tokens a longitud fija."""
    tokens_length = len(tokens)

    # Convertir a tensor 1-D
    tokens = torch.from_numpy(np.array(tokens)).view(-1).contiguous().float()

    if tokens_length < tokens_length_max:
        # Caso corto: añadir zeros de padding al final
        zero_padding = torch.zeros(tokens_length_max - tokens_length)
        tokens = torch.cat([tokens, zero_padding], 0)
    else:
        # Caso largo: truncar a tokens_length_max
        tokens = tokens[:tokens_length_max]
        tokens_length = tokens_length_max

    return tokens, tokens_length
\end{minted}

\section{Flujo de tensores}

\begin{minted}[language=Python, basicstyle=\ttfamily\small, frame=single]{python}
# Entrada
text_ids:     (B, 100)    LongTensor, padded con 0s
text_length:  (B,)        longitudes reales (5-50 típico)

# 1. Embedding
out = embedding(text_ids)
# (B, 100) -> (B, 100, 512)

# 2. Dropout (solo training)
out = dropout(out)
# (B, 100, 512)

# 3. BiGRU con pack_padded_sequence
out = do_RNN_computation(out, text_length, self.gru)
# (B, 100, 512) -> (B, 100, 2048) -> (B, 100, 1024) -> (B, 1024, 100, 1)

# 4. Max-pool sobre tokens (en text_embedding)
out = max(out, dim=2)
# (B, 1024, 1, 1)

# 5. DSC compartido
out = depthwise_seperable_convolution(out)
# (B, 1024, 1, 1)

# 6. squeeze
out = out.squeeze(-1).contiguous()
# (B, 1024, 1)
\end{minted}

\section{Parámetros}

\begin{table}[H]
\centering
\caption{Parámetros de la red de lenguaje.}
\label{tab:lang-params-linea}
\begin{tabular}{lr}
\toprule
\textbf{Componente} & \textbf{Parámetros} \\
\midrule
\texttt{Embedding(29610, 512)} & 15,16 M \\
\texttt{Dropout(0.30)}         & 0 \\
\texttt{BiGRU(512→1024, bidir)} & 4,72 M \\
\texttt{DSC(1024→1024)}        & 10 K \\
\midrule
\textbf{Total lenguaje}        & \textbf{19,89 M} \\
\bottomrule
\end{tabular}
\end{table}
